\capituloSinNumero{Objeto de estudio y campo de acción}
\label{cap:objeto}

\section*{Objeto de estudio}

El objeto de estudio es el sistema de información transaccional (TPS) que
sustenta la operación comercial de Vizcacha Tours. Este tipo de sistema
registra y procesa las transacciones diarias de consulta, cotización, reserva
y pago, y organiza los datos de paquetes, salidas, clientes y pasajeros para
que el personal disponga de información actualizada. El sistema atiende a dos
grupos de usuarios: el personal de la agencia, que administra la información y
confirma las reservas, y los clientes, que acceden con una cuenta propia para
consultar la disponibilidad, solicitar reservas y registrar sus pagos. La
participación de los clientes no genera un proceso paralelo, porque sus
solicitudes ingresan al mismo flujo de validación y de registro que las
atendidas por el personal; esto permite aplicar un criterio uniforme de
disponibilidad y de confirmación. Se trata de un sistema de información
operativo orientado a la gestión comercial, no de un sistema gerencial ni de
un CRM independiente, aunque incluya funciones de registro y seguimiento de
clientes.

\section*{Campo de acción}

El campo de acción se concentra en el control de disponibilidad y la
confirmación de reservas dentro del proceso de venta de paquetes turísticos de
Vizcacha Tours, Potosí. Abarca la administración de paquetes, salidas,
clientes y pasajeros, así como el registro de cotizaciones y pagos, sin
extenderse a la contabilidad general ni a los servicios de otras empresas. La
interacción del cliente se limita a la consulta, a la solicitud de reserva y al
seguimiento de su propia reserva.
